% GENERATED FILE. Edit canonical lessons, then run npm run generate:books.
% canonical-source-hash: fnv1a64:4c9f2145a312852f
% canonical-lessons: LA-C410-rogatio, LA-C410-promulgo, LA-C410-abrogo, LA-C410-sancio, LA-C410-praescribo

\chapter{Passing a Law}
\label{ch:la-passing-a-law}
\hlchaptermodality{410}

\begin{chapteropening}
\textbf{By the end of this chapter:} \emph{I can talk about a bill before the people, publishing and repealing a law, making it binding and laying down what must be done.}

Complete the last lesson of chapter 410: I can talk about a bill before the people, publishing and repealing a law, making it binding and laying down what must be done.
\end{chapteropening}

\section[rogātiō, rogātiōnis]{rogātiō, rogātiōnis — a bill (put to the people); a request}

\label{lesson:LA-C410-rogatio}



\begin{quote}
Before the new one: say the Latin for lawful, then the Latin for allowed.
\end{quote}

\subsection*{You'll want to know: rogātiō, rogātiōnis}
\textbf{rogātiō, rogātiōnis} — \textquotedblleft{}a bill (put to the people); a request\textquotedblright{}.

\begin{quote}
Tribūnus rogātiōnem ad populum tulit.
\end{quote}

\begin{quote}
\emph{The tribune put a bill to the people.}
\end{quote}

\emph{rogātiō, rogātiōnis} f., built on \emph{rogō}, to ask: the magistrate asked the people whether they wished it. If they voted yes, the bill became a \emph{lēx}.

\begin{cousinweb}
\emph{rogātiō} is the noun built on \emph{rogō} \textquotedblleft{}to ask\textquotedblright{} (EtymDB; Lewis \& Short): the question put to the people. English \emph{rogation} comes from it.
\end{cousinweb}

\begin{grammarlens}[title={What you've built}]
One more word for telling, explaining and describing at B1.
\end{grammarlens}

\subsection*{Your turn}
\begin{itemize}
\raggedright
  \item \emph{Say it:} \emph{rogātiō, rogātiōnis}
  \item \emph{Say it:} \emph{rogātiō, rogātiōnis}, once more
  \item \emph{Say it:} a sentence of your own with \emph{rogātiō, rogātiōnis}
  \item \emph{Recall:} say the Latin for lawful, then the Latin for allowed, then say \emph{rogātiō, rogātiōnis} again
\end{itemize}

\begin{tcolorbox}[breakable,colback=teal!4,colframe=teal!35!black,title={Before you move on}]
What does rogātiō, rogātiōnis mean? (\textquotedblleft{}A bill (put to the people); a request\textquotedblright{}.) Say it once more.
\end{tcolorbox}
\section[prōmulgō, prōmulgāre]{prōmulgō, prōmulgāre — to publish, to make known (a bill)}

\label{lesson:LA-C410-promulgo}



\begin{quote}
Before the new one: say the Latin for allowed, then the Latin for a bill.
\end{quote}

\subsection*{You'll want to know: prōmulgō, prōmulgāre}
\textbf{prōmulgō, prōmulgāre} — \textquotedblleft{}to publish, to make known (a bill)\textquotedblright{}.

\begin{quote}
Lēgem novam prōmulgāvit.
\end{quote}

\begin{quote}
\emph{He published the new law.}
\end{quote}

\emph{prōmulgō, prōmulgāre, prōmulgāvī, prōmulgātum}. A bill was posted for public view before the vote. English \emph{promulgate}.

\begin{cousinweb}
The origin of \emph{prōmulgō} is not settled: Lewis \& Short call it unknown. English \emph{promulgate} comes from it.
\end{cousinweb}

\begin{grammarlens}[title={What you've built}]
One more word for telling, explaining and describing at B1.
\end{grammarlens}

\subsection*{Your turn}
\begin{itemize}
\raggedright
  \item \emph{Say it:} \emph{prōmulgō, prōmulgāre}
  \item \emph{Say it:} \emph{prōmulgō, prōmulgāre}, once more
  \item \emph{Say it:} a sentence of your own with \emph{prōmulgō, prōmulgāre}
  \item \emph{Recall:} say the Latin for allowed, then the Latin for a bill, then say \emph{prōmulgō, prōmulgāre} again
\end{itemize}

\begin{tcolorbox}[breakable,colback=teal!4,colframe=teal!35!black,title={Before you move on}]
What does prōmulgō, prōmulgāre mean? (\textquotedblleft{}To publish, to make known (a bill)\textquotedblright{}.) Say it once more.
\end{tcolorbox}
\section[abrogō, abrogāre]{abrogō, abrogāre — to repeal, to abolish}

\label{lesson:LA-C410-abrogo}



\begin{quote}
Before the new one: say the Latin for a bill, then the Latin for to publish.
\end{quote}

\subsection*{You'll want to know: abrogō, abrogāre}
\textbf{abrogō, abrogāre} — \textquotedblleft{}to repeal, to abolish\textquotedblright{}.

\begin{quote}
Populus lēgem veterem abrogāvit.
\end{quote}

\begin{quote}
\emph{The people repealed the old law.}
\end{quote}

\emph{abrogō, abrogāre, abrogāvī, abrogātum}, \emph{ab} + \emph{rogō}: to ask the people to take a law away. English \emph{abrogate}.

\begin{cousinweb}
\emph{abrogō} is built from \emph{ab} \textquotedblleft{}away\textquotedblright{} and \emph{rogō} \textquotedblleft{}to ask\textquotedblright{} (EtymDB; Lewis \& Short): to ask the people to take a law away. English \emph{abrogate} comes from it.
\end{cousinweb}

\begin{grammarlens}[title={What you've built}]
One more word for telling, explaining and describing at B1.
\end{grammarlens}

\subsection*{Your turn}
\begin{itemize}
\raggedright
  \item \emph{Say it:} \emph{abrogō, abrogāre}
  \item \emph{Say it:} \emph{abrogō, abrogāre}, once more
  \item \emph{Say it:} a sentence of your own with \emph{abrogō, abrogāre}
  \item \emph{Recall:} say the Latin for a bill, then the Latin for to publish, then say \emph{abrogō, abrogāre} again
\end{itemize}

\begin{tcolorbox}[breakable,colback=teal!4,colframe=teal!35!black,title={Before you move on}]
What does abrogō, abrogāre mean? (\textquotedblleft{}To repeal, to abolish\textquotedblright{}.) Say it once more.
\end{tcolorbox}
\section[sanciō, sancīre]{sanciō, sancīre — to make sacred and binding; to ratify}

\label{lesson:LA-C410-sancio}



\begin{quote}
Before the new one: say the Latin for to publish, then the Latin for to repeal.
\end{quote}

\subsection*{You'll want to know: sanciō, sancīre}
\textbf{sanciō, sancīre} — \textquotedblleft{}to make sacred and binding; to ratify\textquotedblright{}.

\begin{quote}
Lēgem iūre iūrandō sānxērunt.
\end{quote}

\begin{quote}
\emph{They made the law binding by an oath.}
\end{quote}

\emph{sanciō, sancīre, sānxī, sānctum}: to make a thing sacred and binding. \emph{Sānctus}, holy, is its participle. English \emph{sanction} and \emph{saint}.

\begin{cousinweb}
\emph{sanciō} is inherited from a Proto-Indo-European root meaning \textquotedblleft{}to make holy\textquotedblright{} (EtymDB), and \emph{sānctus} \textquotedblleft{}holy\textquotedblright{} is its participle (Lewis \& Short). English \emph{sanction} and \emph{saint} come from it.
\end{cousinweb}

\begin{grammarlens}[title={What you've built}]
One more word for telling, explaining and describing at B1.
\end{grammarlens}

\subsection*{Your turn}
\begin{itemize}
\raggedright
  \item \emph{Say it:} \emph{sanciō, sancīre}
  \item \emph{Say it:} \emph{sanciō, sancīre}, once more
  \item \emph{Say it:} a sentence of your own with \emph{sanciō, sancīre}
  \item \emph{Recall:} say the Latin for to publish, then the Latin for to repeal, then say \emph{sanciō, sancīre} again
\end{itemize}

\begin{tcolorbox}[breakable,colback=teal!4,colframe=teal!35!black,title={Before you move on}]
What does sanciō, sancīre mean? (\textquotedblleft{}To make sacred and binding; to ratify\textquotedblright{}.) Say it once more.
\end{tcolorbox}
\section[praescrībō, praescrībere]{praescrībō, praescrībere — to lay down, to prescribe}

\label{lesson:LA-C410-praescribo}



\begin{quote}
Before the new one: say the Latin for to repeal, then the Latin for to make sacred and binding.
\end{quote}

\subsection*{You'll want to know: praescrībō, praescrībere}
\textbf{praescrībō, praescrībere} — \textquotedblleft{}to lay down, to prescribe\textquotedblright{}.

\begin{quote}
Lēx praescrībit quid facere dēbeāmus.
\end{quote}

\begin{quote}
\emph{The law lays down what we must do.}
\end{quote}

\emph{praescrībō, praescrībere, praescrīpsī, praescrīptum}: to write at the head, then to lay down a rule. English \emph{prescribe}.

\begin{cousinweb}
\emph{praescrībō} is built from \emph{prae} \textquotedblleft{}before, at the head\textquotedblright{} and \emph{scrībō} \textquotedblleft{}to write\textquotedblright{} (Lewis \& Short). English \emph{prescribe} comes from it.
\end{cousinweb}

\begin{grammarlens}[title={What you've built}]
One more word for telling, explaining and describing at B1.
\end{grammarlens}

\subsection*{Your turn}
\begin{itemize}
\raggedright
  \item \emph{Say it:} \emph{praescrībō, praescrībere}
  \item \emph{Say it:} \emph{praescrībō, praescrībere}, once more
  \item \emph{Say it:} a sentence of your own with \emph{praescrībō, praescrībere}
  \item \emph{Recall:} say the Latin for to repeal, then the Latin for to make sacred and binding, then say \emph{praescrībō, praescrībere} again
\end{itemize}

\begin{tcolorbox}[breakable,colback=teal!4,colframe=teal!35!black,title={Before you move on}]
What does praescrībō, praescrībere mean? (\textquotedblleft{}To lay down, to prescribe\textquotedblright{}.) Say it once more.
\end{tcolorbox}
